\title{Direct Preference Optimization: Your Language Model is Secretly a Reward Model}

\begin{abstract}
Although large unsupervised language models (LMs) acquire extensive world knowledge and exhibit some reasoning capabilities, precise manipulation of their outputs remains challenging due to their training lacking direct supervision. Traditional approaches to improve controllability involve collecting human preference data that ranks model outputs, followed by adapting the unsupervised LM using these preferences—commonly leveraging reinforcement learning from human feedback (RLHF). However, RLHF introduces notable complexity and instability, as it involves first constructing a reward model reflecting the collected preferences, and subsequently fine-tuning the large unsupervised LM using reinforcement learning to optimize this reward while avoiding significant deviation from the pre-trained parameters. In this work, we propose a novel way to parameterize the reward model in RLHF, which allows direct derivation of the optimal policy in closed form. This new perspective transforms the classic RLHF problem into one solvable via a straightforward classification loss. The resulting approach, termed \textit{Direct Preference Optimization} (DPO), offers stability, strong empirical results, and computational efficiency, dispensing with the necessity for generation sampling during fine-tuning or extensive hyperparameter searches. Experimental evaluations reveal that DPO effectively aligns LMs with human preferences on par with or surpassing state-of-the-art approaches. Significantly, DPO not only provides finer control over output sentiment than PPO-based RLHF, but also equals or outperforms existing solutions in areas such as summarization and single-turn conversation—while being much easier to implement and train.
\end{abstract}

\section{Introduction}
Large-scale unsupervised language models (LMs), when trained on extensive datasets, have demonstrated unexpected and advanced capabilities~\citep{chowdhery2022palm, brown2020language, touvron2023llama,bubeck2023sparks}. Nevertheless, such models draw from data created by humans who possess a diverse range of aims, preferences, and expertise. Inevitably, some of these underlying intentions or skillsets are undesirable as behavioral templates. For instance, we may prefer that an AI assistant for programming is able to \textit{recognize} typical coding errors so as to remediate them; at the same time, it is preferable for its code generation to reflect the (sometimes infrequent) instances of high-caliber programming evident in its corpus, rather than mimicking commonplace errors. Likewise, it is useful for a language model to be \textit{informed} about a widespread misconception held by half the population, yet it would be problematic if the model subsequently endorsed this fallacy in 50\% of its relevant outputs. Thus, the essential challenge becomes shaping the model’s \emph{intended behaviors and outputs} from its expansive repertoire of \textit{knowledge and capabilities}, a process crucial to ensuring AI systems that are reliable, effective, and manageable \citep{ouyang2022training}. Whereas prevailing techniques typically guide LMs toward alignment with human values via reinforcement learning (RL), we argue that the objective commonly employed in these RL-based approaches can be precisely recast as a binary cross-entropy objective, thereby offering a significantly more streamlined framework for preference learning.

Broadly, current methodologies impart target behaviors to language models through the use of curated collections of human preferences, which embody norms of safety and helpfulness valued by users. This process of preference learning follows a preliminary stage of large-scale unsupervised pre-training conducted on massive textual corpora. The most direct method for preference learning—supervised fine-tuning on high-quality, human-curated responses—has been overshadowed in efficacy by techniques such as reinforcement learning from human or AI feedback (RLHF/RLAIF; \citep{christiano2017deep,bai2022constitutional}). In RLHF, a reward model is first fit to a dataset containing human preference comparisons, and this model is then leveraged to train the LM policy through RL, favoring outputs rated highly by the reward function, while penalizing excessive deviation from the pre-trained baseline. Despite RLHF’s ability to produce models with notable conversational and programming proficiency, the procedure is far more intricate than standard supervised approaches, necessitating the training of multiple models and frequent on-policy sampling—factors which substantially increase computational demands.

This work demonstrates a method for training a language model to align with human preferences directly, eliminating the need for explicit reward models or reinforcement learning techniques. We introduce \textit{{Direct Preference Optimization} (DPO)}, an algorithm that implicitly targets the same objective as conventional RLHF approaches—maximizing reward under a KL-divergence constraint—but with a much simpler implementation and training process. Conceptually, DPO modifies the log-odds in favor of more preferred responses over less preferred ones, but, crucially, it utilizes a dynamically assigned importance weight for each example. This mechanism counters the mode collapse seen with straightforward probability ratio objectives. DPO, like traditional approaches, builds on theoretical preference models (e.g., the Bradley-Terry model; \cite{bradley1952rankanalysis}) to quantify the agreement between a reward function and observed preference data. Distinctively, instead of using such a model to train a reward function and then deriving a policy to maximize it, DPO applies a variable substitution that allows the preference loss to depend directly on the policy. Provided with human preference data on model outputs, DPO enables direct policy training through a binary cross entropy loss, resulting in an optimal policy for a reward implicitly derived from the preference data.

The principal contribution of this work is Direct Preference Optimization (DPO), a straightforward, reinforcement learning–free approach to preference-based language model training. Experimental results indicate that DPO matches or outperforms existing techniques, including those employing PPO-based RLHF, on a range of preference-driven tasks like sentiment control, summarization, and dialogue, even for models as large as 6B parameters.

Increasingly large self-supervised language models are capable of addressing certain tasks with no additional training (zero-shot; \citep{radford2019language}) or when provided with a handful of example prompts (few-shot; \citep{gpt3,megatron,chowdhery2022palm}). Nevertheless, their outcomes on downstream benchmarks and degree of alignment with user intentions are often significantly enhanced through a process of fine-tuning on corpora composed of task instructions paired with human-crafted responses \citep{mishra-etal-2022-cross,sanh2022multitask,chung2022scaling,thoppilan2022lamda}. This approach, referred to as `instruction-tuning,' fosters LLMs’ capacity to adapt to instructions beyond those seen during training, thereby improving general utility and flexibility \citep{chung2022scaling}. Although instruction tuning is effective, acquiring \textit{relative} human preference annotations is typically less demanding than gathering expert-composed solutions. Consequently, subsequent research has focused on fine-tuning LLMs with datasets that encode human judgments about model outputs, leading to improved results in tasks such as translation \citep{kreutzer-etal-2018-reliability}, summarization \citep{stiennon2022learning,ziegler2020finetuning}, narrative generation \citep{ziegler2020finetuning}, and instruction adherence \citep{ouyang2022training,ramamurthy2023is}. The standard pipeline involves first learning a neural network-based reward function that reflects observed preference data, often under frameworks like the Bradley-Terry model \citep{bradley1952rankanalysis}. The language model is subsequently fine-tuned to optimize this reward, employing reinforcement learning methods such as REINFORCE \citep{williams1992reinforce}, proximal policy optimization (PPO; \cite{schulman2017proximal}), or other variations \citep{ramamurthy2023is}. A related avenue of research utilizes LLMs already tuned for instruction-following, augmented by human feedback, to produce additional synthetic preference examples focused on characteristics like safety or non-harmfulness \citep{bai2022constitutional}. Here, weak human supervision is provided as text-based evaluation criteria used to guide LLM labeling. These techniques blend two significant streams of prior work: the application of reinforcement learning to language models for diverse objectives~\citep{Ranzato2015SequenceLT,paulus2018a,wu2018learning}, and the development of general strategies for learning from human preferences \citep{christiano2017deep,kupcsik2018learning}. Despite the attractiveness of leveraging relative preferences from humans, scaling reinforcement learning-based fine-tuning for large language models still presents considerable technical obstacles; the present study introduces a principled methodology to optimize for relative preferences without depending on RL-based strategies.

Outside the language domain, the problem of learning policies from preference data has been investigated within both bandit and reinforcement learning paradigms, resulting in several distinct methodologies. When preferences or rankings of actions (as opposed to scalar rewards) serve as feedback, the corresponding framework is known as the contextual dueling bandit (CDB; \cite{yue2012karmed,dudik2015contextual}). In scenarios lacking explicit reward signals, the theoretical treatment of CDBs often replaces the standard definition of an optimal policy with the concept of a \textit{von Neumann winner}, i.e., a policy that has an expected probability of at least 50\% to outperform any other policy in head-to-head comparisons \citep{dudik2015contextual}. Despite this, in CDBs, the feedback in the form of preference labels is obtained interactively, whereas in the context of human preference learning, training data is usually derived from a static, finite collection of pre-annotated action pairs \citep{yan2022human}. Along related lines, \textit{preference-based RL} (PbRL) addresses learning from pairwise preferences derived from an unspecified 'scoring' function, substituting direct rewards \citep{BusaFekete2014,ruiz2023dueling}. Multiple PbRL algorithms have been developed, some of which permit the reuse of off-policy preference data; these approaches generally entail first constructing an explicit model of the hidden scoring function (the reward model), followed by its optimization \citep{jain2013learning,BusaFekete2014,christiano2017deep,sadigh2017active,kupcsik2018learning}. By contrast, our method advocates a unified stage in which the policy itself is directly optimized to adhere to the observed preferences.

\section{Preliminaries}\label{section:prelims}

We briefly summarize the RLHF pipeline outlined by \citeauthor{ziegler2020finetuning} and later works \citep{stiennon2022learning, bai2022training, ouyang2022training}, which generally comprises three main steps: (1) supervised fine-tuning (SFT); (2) preference-based reward model training; and (3) reinforcement learning-based policy optimization.

\textbf{SFT}: RLHF commonly initiates by taking a large, pre-trained language model and adapting it to specific target tasks (such as dialogue or summarization) via supervised learning on curated datasets, producing an initial policy $\pi^\text{SFT}$.

\textbf{Reward Modelling Phase}: During the next stage, the SFT-tuned model generates answer pairs $(y_1, y_2)\sim \pi^\text{SFT}(y \mid x)$ for given prompts $x$. Human annotators are then asked to indicate which of the two completions they favor, yielding preference labels expressed as $y_w\succ y_l \mid x$, where $y_w$ and $y_l$ refer to the preferred and non-preferred responses from the sampled pair. It is assumed these preferences are governed by a latent (but inaccessible) reward function $r^*(y, x)$. A variety of statistical models are available to relate observed preferences to latent rewards, with the Bradley-Terry (BT) model \cite{bradley1952rankanalysis} being particularly prominent (although broader ranking frameworks such as the Plackett-Luce model \citep{plackett1975analysis, luce2012individual} can be adopted in situations with more extensive ranking information). Under the BT formulation, the human-derived preference distribution $p^*$ can be expressed as:

Given access to a fixed dataset of pairwise comparisons $\mathcal{D} = \bigl\{x^{(i)}, y_w^{(i)}, y_l^{(i)}\bigr\}_{i=1}^N$ that are drawn from $p^*$, one can model a reward function $r_{\phi}(x, y)$ and determine its parameters through maximum likelihood estimation. When this setup is cast as a binary classification task, the training objective becomes minimizing the following negative log-likelihood:

\begin{equation}\label{eq:reward_model}
    \mathcal{L}_R(r_{\phi}, \mathcal{D}) = -\mathbb{E}_{(x, y_w, y_l)\sim \mathcal{D}}\bigl[\log \sigma(r_{\phi}(x, y_w)- r_{\phi}(x, y_l))\bigr]
\end{equation}

where $\sigma$ denotes the logistic sigmoid. Within the language modeling framework, $r_{\phi}(x, y)$ typically originates from the SFT model $\pi^\text{SFT}(y \mid x)$, extended by a linear output layer atop the last transformer layer, which maps to a scalar reward prediction \cite{ziegler2020finetuning}. Previous studies normalize reward values such that  $\mathbb{E}_{x,y\sim \mathcal{D}}\left[r_\phi(x, y)\right] = 0$ for any context $x$, helping to reduce reward variance.

\textbf{RL Fine-Tuning Phase}: In this phase, the model utilizes the trained reward function to deliver feedback during language model fine-tuning. Following prior work~\citep{jaques2017sequence, jaques2020human}, the optimization problem can be written as:

\begin{equation}\label{eq:RL}
\max_{\pi_{\theta}}  \mathbb{E}_{x\sim \mathcal{D}, y\sim \pi_{\theta}(y \mid x)}\bigl[r_{\phi}(x, y)\bigr] - \beta\mathbb{D}_{\textrm{KL}}\bigl[\pi_{\theta}(y\mid x)\mid \mid \pi_\text{ref}(y\mid x)\bigr],
\end{equation}

Here, $\beta$ governs how far $\pi_\theta$ is permitted to move from the reference policy $\pi_\text{ref}$, typically set to the initial SFT policy $\pi^\text{SFT}$. The language model policy $\pi_\theta$ is commonly initialized as $\pi^\text{SFT}$ as well. Enforcing this regularization is crucial to prevent the policy from drifting too far from regions where the reward model provides reliable estimates and to avoid mode-collapse, thereby preserving output diversity. Because text generation is discrete, the above maximization is non-differentiable and standard practice is to apply reinforcement learning. The prevailing strategy \citep{ziegler2020finetuning, stiennon2022learning, bai2022training, ouyang2022training} defines the reward as ${r(x, y) = r_{\phi}(x, y) -\beta (\log \pi_{\theta}(y\mid x) - \log \pi_\text{ref}(y\mid x))}$ and maximizes it using the PPO algorithm~\cite{schulman2017proximal}.

\section{Direct Preference Optimization}\label{sec:DPO}

Recognizing the considerable complexity and resource demand posed by reinforcement learning when scaling language model fine-tuning, we propose a streamlined alternative that operates directly on preferences. Unlike conventional RLHF pipelines—where a reward function is explicitly learned before being optimized by RL—our method uses a specific reward parameterization permitting analytical derivation of the optimal policy, sidestepping the need for RL altogether. Our core insight is to apply an explicit correspondence between reward functions and optimal policies, which lets us convert an objective over rewards into an objective over policies via a change of variables. This obviates the requirement to explicitly fit a separate reward model, while preserving alignment with established models of human preference such as the Bradley-Terry formulation. Thus, in our approach, the policy network is responsible for modeling both the language generation and the latent reward function.

\textbf{DPO Objective Derivation.} We begin with the standard reinforcement learning objective (Eq.~\ref{eq:RL}) used in earlier studies, which involves an arbitrary reward function $r$. Building upon established literature~\citep{peters2007reinforcement, peng2019advantage, korbak2022reinforcement, go2023aligning}, one can readily demonstrate that the KL-regularized reward maximization objective described in Eq.~\ref{eq:RL} yields the optimal policy of the form:

\begin{equation}\label{eq:op_policy}
    \pi_r(y\mid x) = \frac{1}{Z(x)}\pi_\text{ref}(y\mid x)\exp\left(\frac{1}{\beta}r(x, y)\right),
\end{equation}

with $Z(x) = \sum_{y}\pi_\text{ref}(y\mid x)\exp\left(\frac{1}{\beta}r(x, y)\right)$ acting as the normalization constant, or partition function. A full derivation is available in Appendix~\ref{app:derivation1}. In practice, even with an MLE-based estimate $r_\phi$ of the actual reward function $r^*$, accurately evaluating $Z(x)$ is computationally burdensome~\citep{korbak2022reinforcement, go2023aligning}, rendering this parameterization challenging for practical use. Nonetheless, Eq.~\ref{eq:op_policy} can be re-expressed to isolate the reward as a function of the optimal policy $\pi_r$, the reference policy $\pi_\text{ref}$, and the as-yet-unknown normalization $Z(\cdot)$. Taking the logarithm on both sides and rearranging terms yields:

\begin{equation}\label{eq:main_eq}
    r(x,y) = \beta \log \frac{\pi_r(y\mid x)}{\pi_\text{ref}(y\mid x)} + \beta \log Z(x).
\end{equation}

This relationship holds as well for the ground-truth reward $r^*$ and its corresponding optimal policy $\pi^*$. Conveniently, in the Bradley-Terry model, the likelihood of preferring one completion over another is determined solely by their reward difference: ${p^*(y_1 \succ y_2 \mid x) = \sigma(r^*(x, y_1) - r^*(x, y_2))}$. Substituting the expression from Eq.~\ref{eq:main_eq} for $r^*(x, y)$ into the preference model given in Eq.~\ref{eq:bradley-terry}, we see that the normalization $Z(x)$ cancels out. This allows the human preference probability to be written entirely in terms of the optimal policy $\pi^*$ and the reference policy $\pi_\text{ref}$. Consequently, under the Bradley-Terry model, the optimal RLHF policy $\pi^*$ must satisfy the following:

\begin{equation}\label{eq:objective}
    p^*(y_1\succ y_2 \mid x)=\frac{1}{1 + \exp\left(\beta \log \frac{\pi^*(y_2\mid x)}{\pi_\text{ref}(y_2\mid x)} - \beta \log \frac{\pi^*(y_1\mid x)}{\pi_\text{ref}(y_1\mid x)}\right)}
\end{equation}

Further details of this derivation can be found in Appendix~\ref{app:derivation2}. Although Eq.~\ref{eq:objective} is developed using the Bradley-Terry framework, the same methodology extends to broader Plackett-Luce models~\citep{plackett1975analysis, luce2012individual}, which are explored in Appendix~\ref{app:plackett_luce_models}.

By representing the probability of human preferences using the optimal policy rather than a reward function, we are now positioned to define a maximum likelihood training criterion for any parametric policy $\pi_\theta$. In parallel to the approach of reward model fitting (see Eq.~\ref{eq:reward_model}), our objective for optimizing the policy is:

\begin{equation}\label{eq:optimum_model}
    \mathcal{L}_\text{DPO}(\pi_{\theta}; \pi_\text{ref}) = -\mathbb{E}_{(x, y_w, y_l)\sim \mathcal{D}}\left[\log \sigma \left(\beta \log \frac{\pi_{\theta}(y_w\mid x)}{\pi_\text

In this formulation, we model the implicit reward through an alternative parameterization, resulting in an optimal policy that coincides with $\pi_\theta$. Additionally, because our approach aligns with training a reparameterized Bradley-Terry model, it inherits desirable theoretical guarantees, including consistency under appropriate distributional assumptions for the preference data \cite{bong2022generalized}. We elaborate further on the theoretical attributes of DPO in comparison to related methods in Section~\ref{sec:theory}.

\textbf{How does the DPO update operate?} To gain an operational perspective on DPO, one can examine the gradient of the loss function $\mathcal{L}_\text{DPO}$. With respect to the parameters $\theta$, the gradient takes the form:

\begin{multline*}\label{eq:gradient}
    \nabla_\theta \mathcal{L}_\text{DPO}(\pi_\theta;\pi_\text{ref}) = \\ -\beta\mathbb{E}_{(x, y_w, y_l) \sim \mathcal{D}} \bigg[\underbrace{\sigma(\hat{r}_\theta(x, y_l) - \hat{r}_\theta (x, y_w))}_\text{greater emphasis when the reward estimation is inaccurate}\bigg[\underbrace{\nabla_\theta\log \pi(y_w \mid x)}_\text{promote $y_w$'s probability} - \underbrace{\nabla_\theta\log\pi(y_l \mid x)}_\text{diminish $y_l$'s probability}\bigg]\bigg],
\end{multline*}

where the implicit reward function is given by $\hat{r}_\theta(x, y) = \beta \log \frac{\pi_\theta(y \mid x)}{\pi_\text{ref}(y \mid x)}$, which is defined based on both the model $\pi_\theta$ and the reference $\pi_\text{ref}$ (see further discussion in Section~\ref{sec:theory}). Conceptually, the loss function's gradient acts to increase the log-likelihood of the preferred output $y_w$ while suppressing that of the less favored $y_l$. Critically, the influence of each data point is weighted according to the degree by which the implicit reward model $\hat{r}_\theta$ erroneously rates the less preferred completion higher, modulated by $\beta$. This reflects how severely the implicit reward function misorders outputs, thereby taking the strength of the KL penalty into consideration. Empirical results highlight the crucial role of this weighting mechanism—omitting this factor (i.e., using an unweighted variant) leads to deterioration of the model's behavior, as shown in Appendix Table~\ref{tab:unlikelihood_generations}.

\textbf{DPO overview.} The standard workflow for DPO involves two main steps: First, completions $y_1, y_2 \sim \pi_\text{ref}(\cdot \mid x)$ are generated for each prompt $x$, and these are labeled according to human preferences to form an offline preference dataset $\mathcal{D} = \{x^{(i)}, y_w^{(i)}, y_l^{(i)}\}_{i=1}^N$. Second, the language model $\pi_\theta$ is trained to minimize $\mathcal{L}_\text{DPO}$, conditioned on $\pi_\text{ref}$, the dataset $\mathcal{D}$, and the chosen value of $\beta$.

In practical settings, rather than creating new samples and collecting human annotations, practitioners generally prefer to make use of existing public preference datasets. When such datasets are created with $\pi^\text{SFT}$ as the sampling distribution, we set the initial reference policy as $\pi_\text{ref} = \pi^\text{SFT}$ if it is accessible. In scenarios where $\pi^\text{SFT}$ is not provided, $\pi_\text{ref}$ is established by maximizing the likelihood of the preferred responses ${(x, y_w)}$: specifically, we define ${\pi_\text{ref} = \argmax_{\pi}\mathbb{E}_{x, y_w \sim \mathcal{D}}\left[\log \pi(y_w \mid x)\right]}$. This approach serves to alleviate the distribution mismatch between the true but inaccessible reference distribution and the $\pi_\text{ref}$ used within DPO. Further information regarding implementation details and the selection of hyperparameters is included in Appendix~\ref{app:implementation}.

\section{Theoretical Analysis of DPO} This section deepens the theoretical understanding of DPO, supplies formal justification, and connects the strengths of DPO with the recognized shortcomings of actor-critic methods commonly applied in RLHF (such as PPO~\cite{schulman2017proximal}).

\label{sec:theory}

\subsection{Your Language Model Is Secretly a Reward Model} DPO achieves the dual goals of circumventing both explicit reward model fitting and reinforcement learning for policy optimization, instead relying on a single maximum likelihood framework. Observe that the objective in Eq.~\ref{eq:main_eq} corresponds to a Bradley-Terry model parameterized by a reward function $r^*(x, y) = \beta \log\frac{\pi^*_\theta(y \mid x)}{\pi_\text{ref}(y \mid x)}$. Training our parametric model $\pi_{\theta}$ under this objective is thus equivalent to optimizing the reward model as expressed in Eq.~\ref{eq:reward_model} with an appropriate change of variables. The remainder of this section is devoted to constructing the theoretical foundation for this parameterization, demonstrating its flexibility in capturing the full range of potential reward models, and establishing that it enables exact identification of the optimal policy. Our exposition begins by introducing an equivalence relation among reward functions.

\begin{definition}
Two reward functions $r(x, y)$ and $r'(x, y)$ are said to be equivalent if and only if there exists a function $f$ such that ${r(x, y) - r'(x, y) = f(x)}$.
\end{definition}

This definition clearly forms an equivalence relation, segmenting the space of reward functions into distinct equivalence classes. We formalize this observation in the following two lemmas:

\begin{lemma}\label{lemma:same_prefrence}
Within the framework of Plackett-Luce, and specifically the Bradley-Terry model, any two reward functions belonging to the same equivalence class result in identical preference distributions.
\end{lemma}

\begin{lemma}\label{lemma:same_policy}
For the constrained RL setting, all reward functions from a single equivalence class yield the same optimal policy.
\end{lemma}

The demonstrations of these statements are direct and are included in Appendix~\ref{app:lemma1}. The first lemma highlights a widely recognized issue regarding identifiability in Plackett-Luce models \cite{plackett1975analysis}. Because of this identifiability limitation, additional normalization or constraints are often imposed to ensure meaningful maximum likelihood estimates, as in Eq.~\ref{eq:reward_model} \cite{bong2022generalized}. The second lemma asserts that each reward function within the same equivalence class results in an identical optimal policy, which implies that for the purposes of optimal policy recovery, any representative reward from the equivalence class is sufficient. The following theorem, whose proof appears in Appendix~\ref{app:thm1}, formalizes this intuition:

\begin{theorem}\label{thm:main}
Assuming mild conditions, any reward class consistent with the Plackett-Luce model (and the Bradley-Terry model specifically) can be equivalently written in the reparameterized form ${r(x, y) = \beta \log \frac{\pi(y\mid x)}{\pi_\text{ref}(y\mid x)}}$ for some model $\pi(y\mid x)$ and a fixed reference $\pi_\text{ref}(y \mid x)$.
\end{theorem}

\begin{sproof}
Let $r(x, y)$ be any reward function, and denote by $\pi_r(y \mid x)$ the associated optimal model, as specified in Eq.~\ref{eq:op_policy}. To show that any reward function in the equivalence class of $r$ can be cast in the proposed form, we introduce the following projection:
\begin{equation}
f(r; \pi_\text{ref}, \beta)(x, y) = r(x, y) - \beta\log\sum_{y}\pi_\text{ref}(y\mid x)\exp\left(\frac{1}{\beta}r(x, y)\right)
\end{equation}
This mapping $f$ acts by adjusting $r(x, y)$ via subtraction of the log partition function (which depends solely on $x$), thereby ensuring that $f(r; \pi_\text{ref}, \beta)(x, y)$ remains within the same equivalence class as $r(x, y)$. When substituting the expression for $r$ in Eq.~\ref{eq:main_eq} (valid for any reward function), the projection $f$ yields $f(r; \pi_\text{ref}, \beta)(x, y) = \beta \log \frac{\pi_r(y\mid x)}{\pi_\text{ref}(y\mid x)}$. Thus, this projection produces a representative reward from the equivalence class of $r$ in the desired reparameterized form, without loss of generality.
\end{sproof}

An alternative interpretation of Theorem~\ref{thm:main} is that it uniquely determines the reward function, within each equivalence class, that is chosen by the DPO reparameterization. Specifically, it identifies the reward function that fulfills the condition:

\begin{equation}\label{eq:lag_p}
     \sum_{y}\underbrace{\pi_\text{ref}(y\mid x)\exp\left(\frac{1}{\beta}r(x, y)\right)}_{=\pi(y\mid x)\text{, using Thm.~\ref{thm:main} reparam.}} = 1,
\end{equation}

That is, $\pi(y\mid x)$ constitutes a properly normalized probability distribution (i.e., non-negative probabilities summing to one). Notably, via Eq.~\ref{eq:op_policy}, it becomes apparent that Eq.~\ref{eq:lag_p} serves as the normalization constant (partition function) for the optimal policy derived from $r(x, y)$. A core contribution of the DPO algorithm is the imposition of specific constraints on the Plackett-Luce (and, in particular, Bradley-Terry) class of preference models, which are generally under-determined. These constraints preserve the set of attainable reward models, while also guaranteeing that the optimal policy (as defined in Eq.~\ref{eq:op_policy}) remains analytically tractable for every prompt $x$.

\subsection{Instability of Actor-Critic Algorithms}
Our proposed framework also lends itself to understanding sources of instability in standard actor-critic methods commonly applied in RLHF pipelines, such as PPO. For this analysis, we focus on the RL fine-tuning phase, described in Section \ref{section:prelims}. There is a direct connection to the control as inference perspective~\cite{levine2018reinforcement} applied to the constrained RL setting described in \ref{eq:RL}. Letting $\pi_{\theta}(y\mid x)$ be the parameterized policy model, our goal is to minimize $\mathbb{D}_{\text{KL}}[\pi_{\theta}(y|x) \mid \mid \pi^*(y\mid x)]$, where $\pi^*$ denotes the optimal policy induced by $r_{\phi}(y, x)$, as described in Eq.~\ref{eq:optimum_model}. With some straightforward manipulations, the resulting optimization objective takes the form:

\begin{equation}\label{eq:AC}
    \max_{\pi_{\theta}}\mathbb{E}_{\pi_{\theta}(y\mid x)}\left[\underbrace{r_{\phi}(x, y) -\beta\log\sum_{y}\pi_\text{ref}(y\mid x)\exp\left(\frac{1}{\beta}r_{\phi}(x, y)\right)}_{f(r_{\phi}, \pi_\text{ref}, \beta)} - \underbrace{\beta\log\frac{\pi_{\theta}(y\mid x)}{\pi_\text{ref}(y\mid x)}}_{\text{KL}}\right]
\end{equation}

The objective addressed here aligns with those targeted in earlier research 
\citep{ziegler2020finetuning, stiennon2022learning, bai2022training, ouyang2022training}, employing the DPO-equivalent reward formulation for the class of rewards $r_{\phi}$. Within this framework, the normalization component in $f(r_{\phi}, \pi_\text{ref}, \beta)$ can be understood as the soft value function corresponding to the reference policy $\pi_\text{ref}$. Although this normalization factor does not alter the location of the optimum, omitting it may result in a high-variance policy gradient, thereby impeding the stability of training. Incorporating a learned value function can address the normalization, yet this solution often introduces additional optimization challenges. As another approach, previous work has normalized rewards using a baseline derived from human completions, functioning effectively as a single-sample Monte Carlo approximation of the normalization. The DPO reparameterization distinguishes itself by producing a reward structure that dispenses with the need for any baseline normalization.

\section{Experiments}
This section provides an empirical assessment of DPO for learning policies directly from preference data. We first conduct controlled experiments in text generation to investigate DPO’s capacity to balance reward maximization against KL-divergence minimization relative to the reference policy, and we compare this balance to that achieved by established methods such as PPO for preference learning. Subsequently, we scale our evaluation to larger neural models and address more challenging RLHF scenarios, such as dialogue and summarization tasks. Our observations indicate that, with minimal hyperparameter adjustment, DPO consistently matches or surpasses robust baselines, including PPO-based RLHF and the approach of selecting the best out of $N$ sampled trajectories according to a learned reward. We begin by outlining our experimental procedures, with further specifics provided in Appendix~\ref{app:exp_details}.

\textbf{Tasks.} In our study, we investigate three distinct open-ended text generation tasks. Across all tasks, models are trained to optimize a policy utilizing a preference dataset $\mathcal{D}=\bigl\{x^{(i)}, y_w^{(i)}, y_l^{(i)}\bigr\}_{i=1}^N$. In the \textbf{controlled sentiment generation} setting, $x$ corresponds to an initial segment of a movie review from the IMDb dataset \cite{maas-EtAl:2011:ACL-HLT2011}, and the objective is for the policy to output $y$ that reflects positive sentiment. For this controlled experiment, we \textit{construct} preference pairs between generations by employing a pre-trained sentiment classifier such that $p(\text{positive}\mid x,y_w)>p(\text{positive}\mid x,y_l)$. The SFT baseline involves fine-tuning GPT-2-large on the training split of IMDb reviews until convergence (additional information provided in App~\ref{app:sentiment_details}). In the \textbf{summarization} task, $x$ represents a Reddit forum post, and the policy must produce a summary $y$ highlighting the core arguments or points. Adhering to established approaches, we use the Reddit TL;DR summarization dataset \citep{volske-etal-2017-tl} and leverage human preference annotations collected by \citeauthor{stiennon2022learning}. Our SFT model for this task is trained on summaries written by humans for forum posts\footnote{\url{https://huggingface.co/CarperAI/openai_summarize_tldr_sft}} via the TRLX \citep{leandro_von_werra_2023_7790115} framework to enable RLHF. Notably, human preference data was acquired from outputs generated by a separate SFT model that was trained similarly. Lastly, for the \textbf{single-turn dialogue} task, $x$ is a user-generated prompt, which could range from scientific inquiries to personal advice requests. The model’s goal is to produce a helpful and engaging response $y$; for this, we utilize the Anthropic Helpful and Harmless dialogue dataset \citep{bai2022training}, comprised of 170,000 exchanges between humans and automated agents. Each dialogue concludes with two candidate responses from a large (though unspecified) language model and a corresponding human preference label identifying the superior response. Since a dedicated SFT model does not exist for this scenario, we fine-tune an available pretrained language model exclusively on the responses labeled as preferred, yielding the SFT baseline.

\textbf{Evaluation.} Our experimental evaluation employs two distinct methodologies. To rigorously assess each algorithm’s capability in maximizing the constrained reward, we examine their performance in controlled sentiment generation. Here, we compute the reward and KL-divergence frontier for each algorithm, leveraging our access to the true reward function—specifically, a sentiment classifier—to obtain this frontier. Outside this controlled setting, in practical scenarios where the true reward is unavailable, we gauge performance by the \textit{win rate} over a baseline policy. For this, we utilize GPT-4 as a stand-in for human judgment, measuring perceived summary quality in summarization and helpfulness in single-turn dialogue tasks. In the summarization experiments, the baseline consists of test set reference summaries; in dialogue, the baseline is the preferred human response from the dataset. While prior work demonstrates that language models may surpass traditional automatic metrics as evaluators \citep{Chen2023ExploringTU}, we supplement our study with a human evaluation to substantiate our use of GPT-4 as an assessor (detailed in Sec.~\ref{sec:human-judgments}). Our findings reveal strong correspondence between GPT-4 and human raters, with humans agreeing with GPT-4 at rates equal to or exceeding their agreement with one another.

\textbf{Methods.} Beyond DPO, we benchmark several established methods for training language models to align with human preferences. For summarization, we test zero-shot prompting with \textbf{GPT-J} \citep{gpt-j}, and for the dialogue setting, we implement two-shot prompting using \textbf{Pythia-2.8B} \citep{biderman2023pythia}. Our evaluation also encompasses the \textbf{SFT} model and \textbf{Preferred-FT}, which involves supervised fine-tuning on preferred completions $y_w$, using outputs from either the SFT model (for sentiment and summarization) or a general LM (for single-turn dialogue). We further assess the pseudo-supervised \textbf{Unlikelihood} approach~\citep{welleck2019neural}, which seeks to increase the likelihood of $y_w$ while explicitly decreasing the likelihood of $y_l$, modulated by an optional coefficient $\alpha\in[0,1]$ for the unlikelihood component. Additionally, we evaluate \textbf{PPO} \citep{schulman2017proximal} where the reward model is trained on preference annotations, alongside \textbf{PPO-GT}, which has access to the gold-standard reward in the controlled sentiment task. For sentiment analysis, two PPO-GT variants are included: a standard off-the-shelf implementation \cite{leandro_von_werra_2023_7790115}, and a tailored version with reward normalization and tuned hyperparameters to enhance performance; these improvements are likewise applied to PPO with learned rewards. As a final comparison, we utilize the \textbf{Best of $N$} baseline: $N$ samples are drawn from the SFT model (or Preferred-FT in dialogue), with the optimal response selected based on a reward model learned from preference data. While this strategy decouples reward estimation from PPO optimization and typically performs strongly, its computational demands grow rapidly with $N$, rendering it inefficient for practical use during evaluation, as it necessitates $N$ samples per query.

\subsection{How well can DPO optimize the RLHF objective?}

The standard reward maximization objective employed in RLHF algorithms, which incorporates a KL divergence constraint, is designed to strike a balance between maximizing the reward signal and ensuring the policy does not diverge excessively from the reference policy. Thus, when assessing the performance of various methods, it is essential to jointly consider both the magnitude of the obtained reward and the level of KL divergence; a higher reward accompanied by a significant increase in KL is not always preferable. Figure~\ref{fig:frontier-tldr-main} illustrates the tradeoff between reward and KL for several approaches within the sentiment classification context. For each method, we conduct a series of training runs, each with a distinct hyperparameter setting to control policy conservativeness (target KL values of $\{3,6,9,12\}$ for PPO; $\beta$ values of $\{0.05,0.1,1,5\}$; and $\alpha$ values of $\{0.05,0.1,0.5,1\}$ for unlikelihood; and different random seeds for preferred-FT), yielding a total of 22 training configurations. At every 100 training iterations until the algorithms converge, we assess each resulting policy on a hold-out set of prompts, measuring the average reward given the ground-truth reward function and the mean sequence-level KL\footnote{Here, this refers to the aggregate sum of KL divergences at each timestep.} relative to the reference policy, specifically $\text{KL}\left(\pi\mid \mid \pi_\text{ref}\right)$. Our analysis reveals that DPO achieves the most favorable tradeoff frontier, attaining superior reward outcomes at comparably low KL values. This observation is particularly striking for several reasons: despite DPO and PPO optimizing the same underlying objective, DPO consistently achieves greater efficiency, strictly dominating PPO in terms of the reward versus KL tradeoff. Moreover, DPO's performance remains superior even when PPO is provided access to exact reward signals (PPO-GT).

\subsection{Can DPO scale to real preference datasets?}
\label{sec:dpo-real-datasets}
We proceed to assess how well DPO fine-tunes on real-world preference datasets in the contexts of summarization and single-turn dialogue. In summarization tasks, it has been observed that standard automatic evaluation metrics like ROUGE often exhibit weak alignment with human judgments~\citep{stiennon2022learning}. Previous studies have demonstrated that fine-tuning language models via PPO using human preference data can yield summaries that are more favorably received by humans. Our evaluation framework involves generating completions on the TL;DR summarization dataset's test set, then determining each method’s average win rate when compared to the test set’s reference completions. Completions are sampled from all models with temperatures spanning 0.0 to 1.0; the resulting win rates are presented in Figure~\ref{fig:frontier-tldr-main} (right). DPO, PPO, and Preferred-FT all fine-tune the same initial GPT-J SFT model\footnote{\url{https://huggingface.co/CarperAI/openai_summarize_tldr_sft}}. Our experiments reveal that DPO attains a win rate near 61\% at a temperature of 0.0, outperforming PPO, which reaches about 57\% at its optimal setting of 0.0. In addition, DPO surpasses the best of $N$ baseline in terms of peak win rate. It is important to mention that the DPO $\beta$ hyperparameter was not substantially tuned in our experiments, so the actual potential of DPO may be higher than reflected here. Additionally, DPO demonstrates greater robustness to changes in sampling temperature compared to PPO, whose performance may decline to match that of the baseline GPT-J at elevated temperatures. Notably, Preferred-FT does not deliver meaningful improvements over the original SFT model. Furthermore, in human evaluations detailed in Section~\ref{sec:human-judgments}, direct comparison reveals that samples from DPO at temperature 0.25 were favored 58\% of the time over those from PPO at the same temperature.

For the single-turn dialogue evaluation, we focus on examples from the test set of the Anthropic HH dataset \citep{bai2022training} involving just one exchange between a human and an assistant. In assessments using GPT-4, we employ the test set's preferred completions as a benchmark to calculate win rates across different approaches. Since a standard SFT model is not established for this task, our pipeline initiates from a pre-trained Pythia-2.8B model; we fine-tune this model with Preferred-FT on the selected completions to create a reference model that produces in-distribution completions, followed by training with DPO. To contextualize DPO's results, we compare it with the strongest among 128 Preferred-FT completions (as our analysis indicates that the Best of $N$ metric saturates at $N=128$ for this scenario; refer to Appendix Figure~\ref{fig:best-of-n}) as well as a 2-shot prompt configuration of the unmodified Pythia-2.8B model, observing that DPO achieves equivalent or superior win rates at optimal temperatures for each respective model. We further assess an RLHF model trained via PPO on the Anthropic HH dataset \footnote{\url{https://huggingface.co/reciprocate/ppo_hh_pythia-6B}} provided by a recognized source \footnote{\url{https://github.com/CarperAI/trlx/tree/main/examples/hh}}, but despite varying prompts and sampling temperatures, we do not observe any improvements over the original Pythia-2.8B. Drawing on findings from TL;DR and given that both Best of 128 and PPO maximize the same reward signal, we use Best of 128 as an approximate indicator of PPO performance. Notably, DPO emerges as the only method that efficiently exceeds the quality of the preferred completions on the Anthropic HH dataset, delivering performance on par with or better than the computationally intensive Best of 128 approach. Figure~\ref{fig:dialogue-main} further demonstrates that DPO reaches peak effectiveness in relatively few training steps.

\subsection{Generalization to a new input distribution}

To further assess the robustness of PPO and DPO when exposed to distributional changes, we test the policies trained in our Reddit TL;DR summarization setup on a distinct dataset: the test split of the CNN/DailyMail news articles \citep{nallapati-etal-2016-abstractive}. For this evaluation, we apply the optimal sampling temperatures identified from TL;DR (0 and 0.25). The outcomes, shown in Table~\ref{tab:ood}, are determined using GPT-4 to rate the models' outputs against ground-truth summaries. The evaluation used the same GPT-4 (C) prompt as in Reddit TL;DR, with the term “forum post” switched to “news article”. Across this new dataset, DPO maintains a clear advantage over PPO. These findings offer preliminary evidence that DPO can generalize to new input distributions on par with, or better than, PPO, despite not leveraging the extra unlabeled Reddit TL;DR prompts employed by PPO.

\subsection{Validating GPT-4 judgments with human judgments}
\label{sec:human-judgments}
To determine how well GPT-4’s assessments align with human preferences, we carry out a human evaluation on outputs from the TL;DR summarization task, utilizing two separate GPT-4 prompts. The \textbf{GPT-4 (S)} (simple) prompt asks for a selection of the summary that best covers the essential information from the post. In contrast, the \textbf{GPT-4 (C)} (concise) prompt requests the more concise and informative summary, addressing our observation that the \textbf{GPT-4 (S)} prompt tends to prefer verbose or repetitive summaries, differing from human preferences. Detailed versions of the prompts are available in Appendix~\ref{app:prompts}. We analyze three scenarios by contrasting the top-performing approach (DPO, temp. 0.25), the lowest-performing one (PPO, temp. 1.0), and an intermediate method (SFT, temp. 0.25) against greedily-sampled PPO at its optimal temperature; this set was chosen to span a range of sample qualities. Our analysis reveals that, for both prompts, GPT-4’s decisions agree with human judgments at rates comparable to inter-human agreement, indicating GPT-4’s validity as a stand-in for human evaluation (note that we gathered multiple human ratings only for DPO and PPO-1 comparisons due to rater availability). The \textbf{GPT-4 (C)} prompt, in particular, yields win rates that closely reflect human preferences and is thus adopted for reporting key findings in Section~\ref{sec:dpo-real-datasets}. Additional specifics about the human rating study—including screenshots of the user interface and information about the participating volunteers—can be found in Appendix~\ref{app:human-study}.

\section{Discussion}
Preference-based learning represents an effective and extensible approach for developing robust and aligned language models. In this work, we have presented DPO, a streamlined methodology for language model training from preference data that circumvents reinforcement learning entirely. Rather than forcing the preference alignment challenge into a typical reinforcement learning framework and leveraging conventional RL tools, DPO establishes a direct correspondence between policy optimization in language models and associated reward functions. This allows for direct alignment to human preferences using a straightforward cross-entropy loss, bypassing both reinforcement learning and any compromise in representational capability. Remarkably, DPO matches or surpasses the performance of established RLHF strategies, such as those built on PPO, often without extensive hyperparameter adjustments. Consequently, DPO significantly lowers the practical challenges associated with preference-driven language model development.

\textbf{Limitations \& Future Work.} This study highlights several directions that warrant further investigation. An open question is how the DPO-trained policies behave on distributions distinct from those seen during training, particularly in comparison to policies trained using explicit reward functions. Preliminary evidence points to comparable generalization capabilities between DPO and PPO-based approaches; however, systematic analysis is necessary. For instance, it remains to be explored whether the self-labeling strategy using DPO-generated outputs can capitalize on unlabeled data in an efficient manner. Moreover, understanding the occurrence and impact of reward over-optimization in the context of direct preference optimization is crucial; it is unclear if the minor drop observed in performance in Figure~\ref{fig:dialogue-main}-right is attributable to this phenomenon. Although our experiments span models containing up to 6B parameters, extending DPO training to much larger, state-of-the-art architectures poses a compelling opportunity. With respect to evaluation, our findings suggest that GPT-4’s calculated win rates are sensitive to prompt formulation; identifying optimal prompting strategies for eliciting high-quality assessments from automated evaluators is a valuable avenue for future study. Ultimately, the applicability of DPO extends well beyond language modeling from human preferences and could potentially benefit generative modeling tasks in a variety of domains.